\documentclass{article}
\usepackage{graphicx}

\begin{document}

\section{引言}
\label{sec:intro}

本文设计了一个用于论文文本处理的系统，并完成了相关模块的实现。该系统通过任务状态记录、段落抽取和模型改写流程，提供了一种较为完整的处理方案，具有重要意义。

\section{方法}
\label{sec:method}

如图 \ref{fig:arch} 所示，系统首先读取用户选择的 LaTeX 项目，然后解析报告中的文本片段，最后对需要处理的段落生成修订建议。该过程保留原有引用和交叉引用，例如 \cite{smith2024}。

\begin{figure}
\centering
\fbox{architecture}
\caption{System architecture}
\label{fig:arch}
\end{figure}

\bibliographystyle{plain}
\bibliography{refs}

\end{document}
